%% kapa: exam=1 round=37 year=2026 period=1 form=A subject=real_estate subject_ko=부동산학원론 seq=13 number=93
%% source: Q-Net 감정평가사 기출문제 (cst00309) artlSeq=5258924 file=raw/1008/5258924/2026년 제37회 감정평가사 1차 1교시.pdf
%% conversion: kordoc 4.24.51 -> tools/gen1.py (auto); review=proofread-sonnet; flags=none
\begin{question}{93}
효율적 시장이론(efficient market theory)에 관한 설명으로 옳은 것은?
\begin{choices}
  \item 효율적 할당(efficient allocation)이란 위험과 무관하게 모든 투자대안의 수익률이 동일하도록 자원이 할당되는 것을 말한다.
  \item 강성 효율적 시장은 완전경쟁시장의 가정에 부합되는 시장이며, 어느 시장참여자도 초과이윤을 획득할 수 없다.
  \item 약성 효율적 시장은 공표된 모든 정보를 지체없이 시장가치에 반영하는 시장이므로 기술적 분석을 통해 초과이윤을 획득할 수 없다.
  \item 약성ㆍ준강성ㆍ강성 효율적 시장 모두 정보비용이 존재하지 않는다.
  \item 불완전경쟁시장은 할당 효율적 시장이 될 수 없다.
\end{choices}
\end{question}
\answer{2}
